
\documentclass[12pt]{article}

% Packages
%-----------------------------------------------------------------

% Allow direct use of accents such as á é ñ.
\usepackage[utf8]{inputenc}

% This is a good idea to have some symbols included within the font
% properly:
% http://tex.stackexchange.com/questions/664/why-should-i-use-usepackaget1fontenc
\usepackage[T1]{fontenc}


% Set type of paper and margins.
\usepackage[a4paper, margin=2.7cm, % Set the size and margins of the paper
includehead,
nomarginpar,% We don't want any margin paragraphs
headheight=10mm,% Set \headheight to 10mm
]{geometry}




\usepackage{amsmath, amsthm, amsfonts, amssymb}

% Generate a PDF with hyperlinks in references.
\usepackage[colorlinks=true,linkcolor=black,citecolor=black,urlcolor=black,breaklinks]{hyperref}



\usepackage{graphicx, xcolor}
\usepackage{placeins}
\usepackage{fancyhdr}
% Source - https://tex.stackexchange.com/a/13407
% Posted by lockstep, modified by community. See post 'Timeline' for change history
% Retrieved 2026-03-11, License - CC BY-SA 3.0

\renewcommand{\footrulewidth}{0.4pt}% default is 0pt

%Styling
\setlength{\parindent}{0em} 


% Uncomment for line numbers
% \usepackage{lineno}
% \linenumbers

% Bibliography
%-----------------------------------------------------------------

% This uses a bibliography style which hyperlinks the paper titles to
% the paper URL specified in the bibtex file. It also uses natbib,
% which cites papers by name such as Euler (1770) instead of [17].


\usepackage{natbib}
\usepackage{url}
\bibliographystyle{apalike}
%\bibliographystyle{plainnat-linked}
% \bibliographystyle{plain}



% Shortcuts
%-----------------------------------------------------------------
% none defined at the moment.


% Title, author, date
%-----------------------------------------------------------------

% If set, these will be the internal title and author of the PDF (and
% will be listed, for example, in ereaders and tablets.

\hypersetup{pdftitle={ProofMonthDayYearIsWrong}}
\hypersetup{pdfauthor={Brueckmann}}


% Paper title and author
\title{Why the Month Day Format is not only absurd but also a wrong date format -- a proof by contradiction}
\author{\href{https://brueckmann.github.io/}{G. Brückmann}}

% Date is set automatically unless specified below.
 \date{March 11, 2026}
 %It pains me to write it like this. 


% Addresses. The "article" class does not have a standard way to
% include them. This imitates the behaviour of "amsart" by inserting
% them at the end of the paper.
%-----------------------------------------------------------------

% \AtEndDocument{\bigskip{\footnotesize
%   \noindent
%   \textsc{G. Brückmann. Please do not contact me because of this.}
%   % \emph{E-mail address}: \texttt{\href{mailto:gracia.brueckmann@unibe.ch}{gracia.brueckmann@unibe.ch}}
% }}

% =========================================================================

\begin{document}

% Set the page style to "fancy"...
\pagestyle{fancy}
%... then configure it.
\fancyhead{} % clear all header fields
\lhead{G. Brückmann}
\rhead{\textbf{Month Day is wrong}}
\lfoot{\href{https://raw.githubusercontent.com/brueckmann/Proof_MonthDayYearIsWrong/refs/heads/main/Proof_MonthDayYearIsWrong.pdf}{Link to the most recent version of this article:} \href{https://raw.githubusercontent.com/brueckmann/Proof_MonthDayYearIsWrong/refs/heads/main/Proof_MonthDayYearIsWrong.pdf}{\url{github/brueckmann/Proof_MonthDayYearIsWrong/blob/main/Proof_MonthDayYearIsWrong.pdf}}}
\rfoot{\thepage}


\fancypagestyle{plain}{%
  \fancyhf{}%
  \renewcommand{\footrulewidth}{0.1mm}%
\fancyfoot[L]{\href{https://raw.githubusercontent.com/brueckmann/Proof_MonthDayYearIsWrong/refs/heads/main/Proof_MonthDayYearIsWrong.pdf}{Link to the most recent version of this article:} \href{https://raw.githubusercontent.com/brueckmann/Proof_MonthDayYearIsWrong/refs/heads/main/Proof_MonthDayYearIsWrong.pdf}{\url{github/brueckmann/Proof_MonthDayYearIsWrong/blob/main/Proof_MonthDayYearIsWrong.pdf}}}
  \fancyfoot[R]{\thepage}%
  \renewcommand{\headrulewidth}{0mm}%
} 


\maketitle

\begin{abstract}
The date above is wrong. It is not a correct representation of the actual date, which is 2026-03-11, or the 11$^\text{th}$ of March, 2026. This is why. Therefore, please retire this date format so that we all have good dates.\\
\par
\emph{Keywords: } Worst, date, format, ever, proof
\end{abstract}

% \tableofcontents

\section{Introduction}
Many people ask themselves, \emph{What is the best date}? While answers may vary, as this might be subjective, within this article, the author proves that the date as displayed above (\emph{March 11}) \emph{is wrong}, making it the worst of the common dates. Also called the \emph{`American absurdity'} \citep{redditor}, Month Day is widely recognised as the worst date format. That it is also objectively wrong is proven in this article.

\section{Proof}


Assume 11$^{\text{th}}$ of March, 2026 is today's date, and let's agree this is correct. Let's also consider Mary, who (famously) has a little lamb (and disregard that we don't know how the lamb feels about being possessed by Mary). \\

Mary's little lamb implies the little lamb of Mary. Let's apply something similar to the correct date, the 11$^{\text{th}}$ of March, which, recall, is a true statement.\\ 

\clearpage


\begin{align*} 
\text{little lamb of Mary} &=  \text{Mary's little lamb} \\
\text{Mary's little lamb}  &= \text{little lamb of Mary}\\
11^{\text{th}}  \text{ of March} &=  \text{March's } 11^{\text{th}}\\
\end{align*}
It becomes evident that the 11$^{\text{th}}$ of March is indeed the 11$^{\text{th}}$ (day) of March.\\

Now, take the case of March 11 in order to prove by contradiction and suppose: \\ 
\begin{align*}
11^{\text{th}}  \text{ of March} &= \text{March } 11   \\ 
\text{But as:} \\
\text{little lamb of Mary} &\neq  \text{Mary little lamb}
\end{align*}
\hfill $\blacksquare$

Given \emph{Mary little lamb} is so ambiguous that you may wonder if the author is under the influence, one can confidently refute the claim that March 11 is a valid date. 

\section{Conclusion}
 
We can confidently conclude that March 11 is not equivalent to today's date, which is 2026-03-11, the best date format IMHO, and \citeauthor{ISO8601} agrees with me. Given the proof in this paper, we should retire the Month Day Year date format, once and for all.



\bibliography{library}





\makeatletter\let\ps@plain\ps@fancy\makeatother % anywhere after the first page/chapter

% \vspace*{\fill}
% \footnotesize{
% \href{https://raw.githubusercontent.com/brueckmann/Proof_MonthDayYearIsWrong/refs/heads/main/Proof_MonthDayYearIsWrong.pdf}{Link to the most recent version.} \\
% \href{https://raw.githubusercontent.com/brueckmann/Proof_MonthDayYearIsWrong/refs/heads/main/Proof_MonthDayYearIsWrong.pdf}{
% \url{https://raw.githubusercontent.com/brueckmann/Proof_MonthDayYearIsWrong/refs/heads/main/Proof_MonthDayYearIsWrong.pdf}
% }


\end{document}

% quelle: Vorlage https://canizo.org/templates#article_class